\documentclass[11pt]{article}
\usepackage[a4paper,margin=25mm]{geometry}
\usepackage{amsmath,amssymb,amsthm}
\usepackage[hidelinks]{hyperref}
\newtheorem{theorem}{Theorem}
\newcommand{\E}{\mathbb E}
\newcommand{\Prob}{\mathbb P}
\title{Zero mean does not imply a symmetric\\rough L\'evy-area sign law\\
\large Disproof of Conjecture 00000004027}
\author{gaochengzhecpu}
\date{}
\begin{document}
\maketitle
\begin{abstract}
A planar pure-area rough path has zero first level and second-level increments
$\mathbb X_{s,t}=A(t-s)J$, where $J$ is the standard antisymmetric matrix and
$A$ takes the values $2,-1,-1$ on three equiprobable outcomes.
Every entry of $\E\mathbb X_{s,t}$ vanishes for every time interval.
Nevertheless, the area sign law is constantly
$\frac13\delta_1+\frac23\delta_{-1}$.
Reflection leaves the entire mean tensor unchanged and reverses the sign bias.
This disproves both the asserted sufficient condition and the claim that the
asymmetry is determined by that mean.
\end{abstract}

\section*{The assertion and the rough-path convention}
The source claims that, for rough planar drivers, the area sign distribution
converges to a symmetric law if and only if the second level has zero mean,
and that its asymmetry is a functional of that mean. No stationarity,
independence, mixing, or nondegeneracy hypothesis is stated. The example
below satisfies zero mean for the \emph{whole tensor at all times}, not merely
for its antisymmetric part or at selected endpoints.

We use the standard weak geometric step-2 rough-path convention.
Its increments $(X_{s,t},\mathbb X_{s,t})$ satisfy
\begin{align*}
X_{s,t}&=X_{s,u}+X_{u,t},\\
\mathbb X_{s,t}&=\mathbb X_{s,u}+\mathbb X_{u,t}
                  +X_{s,u}\otimes X_{u,t},\\
\mathbb X_{s,t}^{ij}+\mathbb X_{s,t}^{ji}&=X_{s,t}^iX_{s,t}^j,
\end{align*}
with bounds $|X_{s,t}|\leq C|t-s|^{1/2}$ and
$|\mathbb X_{s,t}|\leq C|t-s|$ on compact time intervals.
We use maximum coordinate norms; finite-dimensional norm equivalence makes
this choice immaterial. These bounds also imply every weaker rough-path
H\"older regularity $1/3<\alpha<1/2$ locally.

Pure-area paths are standard rough enhancements of the zero first-level
path. For the terminology, see Caravenna, Gubinelli, and Zambotti,
\emph{Lectures on rough paths}, version 14 May 2025,
Section 11.4, Lemma 11.9: an additive second level with the required regularity
gives a pure-area rough path, and antisymmetry gives weak geometricity.
All defining identities and bounds needed here are checked directly below.

\section*{A genuine pure-area rough driver}
Let $\Omega=\{0,1,2\}$ with probability $1/3$ at each point, and set
\[
A(0)=2,\qquad A(1)=A(2)=-1,
\qquad J=\begin{pmatrix}0&1\\-1&0\end{pmatrix}.
\]
For each outcome define, for every $s,t\in\mathbb R$,
\[
X_{s,t}=0,\qquad \mathbb X_{s,t}=A(t-s)J.
\]
Both levels vanish on the diagonal. Additivity of $t-s$ proves Chen's
identity, since its cross term is zero. Antisymmetry of $J$ proves the shuffle
identity. Entrywise,
\[
|X_{s,t}^i|=0,\qquad
|\mathbb X_{s,t}^{ij}|\leq |A|\,|t-s|.
\]
Thus every outcome is a weak geometric rough path, with the stated regularity;
one can even use the common deterministic constant $C=2$.
The second level is an enhanced rough-path component, rather than the
classical iterated integral of the identically zero smooth path.

\begin{theorem}
The whole second-level mean vanishes on every interval, but its area sign
law is nonsymmetric and remains unchanged on every nondegenerate interval.
\end{theorem}
\begin{proof}
Since $\E A=(2-1-1)/3=0$,
\[
\E\mathbb X_{s,t}^{ij}=(t-s)J_{ij}\E A=0
\quad\text{for all }s,t,i,j.
\]
Use the conventional antisymmetric normalization of L\'evy area,
\[
L_{s,t}=\frac{\mathbb X_{s,t}^{12}-\mathbb X_{s,t}^{21}}2=A(t-s).
\]
For every $s<t$, its sign equals that of $A$ and never vanishes. Hence
\[
\operatorname{Law}(\operatorname{sgn}L_{s,t})
=\frac13\delta_1+\frac23\delta_{-1}.
\]
This law is unequal to its reflection. It is constant, so its limit is the
same nonsymmetric law. In particular, both the long-time interpretation
$L_{0,t}$ as $t\to\infty$ and the small-time interpretation as $t\downarrow0$
fail. The same holds along any sequence of nondegenerate intervals.
A different positive normalization of area does not change the sign.
\end{proof}

\section*{The entire mean cannot determine the asymmetry}
Replace $A$ by $-A$. All rough-path conditions still hold, and the complete
function $(s,t,i,j)\mapsto\E\mathbb X_{s,t}^{ij}$ is still identically zero.
But the reflected area has sign law
$\frac23\delta_1+\frac13\delta_{-1}$.
Consequently the signed asymmetry
\[
b=\Prob(L_{s,t}>0)-\Prob(L_{s,t}<0)
=\E[\operatorname{sgn}L_{s,t}]
\]
equals $-1/3$ for the first driver and $1/3$ for its reflection.
No functional of the entire mean tensor, explicit or otherwise, can recover
this asymmetry for every driver. This also explains the basic obstruction:
zero mean balances magnitudes, whereas symmetric signs balance probabilities.

\section*{Connection with actual smooth signatures}
For completeness, the supplied Lean file also checks the smooth planar loop
\[
\gamma_a(u)=\bigl(u-u^2,\;60a(u^2-u^3)\bigr),\qquad 0\leq u\leq1.
\]
It starts and ends at zero, and direct integration gives its actual second
signature
\[
\int_0^1 \gamma_a^i(u)\,d\gamma_a^j(u)=aJ_{ij}.
\]
The file proves all four entries, using actual derivatives and interval
integrals. It also defines concatenation with the prefix increments and
proves that $N$ repeated copies have second signature $NaJ$.
These loops are supplementary: their interior-time diagonal means need not
vanish, so the all-time counterexample above does not use that property.

The same loops give the usual smooth-lift approximation of the pure-area
example with a loss of regularity. On each cell of a mesh of size $1/n$,
traverse $n^{-1/2}\gamma_a$, rescaled in time. The first-level excursion is
$O(n^{-1/2})$, and each completed cell contributes exactly $aJ/n$.
On compact intervals the resulting lifts have uniform $1/2$-H\"older bounds
and converge uniformly, in both levels, to $(0,a(t-s)J)$.
Interpolation therefore gives convergence in every weaker
$\alpha$-H\"older rough-path topology, $1/3<\alpha<1/2$.
Piecewise smooth loops may be smoothed without changing the limit.
Thus the example also lies in the geometric class at any such exponent
(and in $p$-variation for $2<p<3$). The approximation and interpolation
argument in this paragraph is a mathematical explanation, not a theorem
formalized in the accompanying Lean project.

\section*{Exact formalization scope}
The primary formalization is in the namespace \texttt{PureArea}.
Its structure \texttt{IsWeakGeometricHalf} spells out the actual two Chen
identities, diagonal conditions, shuffle identity, and H\"older bounds.
The theorem \path{genuine_weak_geometric_rough_path} proves every field;
\path{all_time_whole_tensor_mean_zero} quantifies over arbitrary real
times and both tensor indices. This is a definition-level verification;
it does not claim an existing Mathlib rough-path library.

The probability space is normalized counting measure on \texttt{Fin 3}.
The area is extracted from the second tensor, and its sign law is an actual
measure pushforward. The file proves the complete law is constant, its two
atom masses are $1/3$ and $2/3$, and it differs from its reflected measure.
It further proves that along any sequence of nondegenerate intervals the two
atom probabilities cannot have a common limit. For laws supported on
$\{-1,1\}$ this is exactly the necessary probability condition for a
symmetric limiting law. No general topology of weak convergence of measures
is developed. Finally, the reflection argument formally rules out an
arbitrary functional of the full mean tensor that returns expected sign.

The project uses Lean 4.19.0 and Mathlib at a fixed commit. The printed
dependencies contain only the standard axioms \texttt{propext},
\texttt{Classical.choice}, and \texttt{Quot.sound}.
There are no added axioms, proof placeholders, or numerical oracles.

\section*{Sources}
\begin{enumerate}
\item \href{https://github.com/The-Last-Math-Competition/The-Last-Math-Competition/blob/9b795e7a94a6076e49a65a6489e1caf9153abd23/conjectures/00000004027.md}{The Last Math Competition, Conjecture 00000004027}.
The exact bilingual source bytes and their commit and SHA-256 provenance
are included with the submission.
\item F. Caravenna, M. Gubinelli, and L. Zambotti,
\href{https://www.lpsm.paris/_media/users/zambotti/bookcgz.pdf}{\emph{Lectures on rough paths}},
version 14 May 2025, Section 11.4, Lemma 11.9.
\end{enumerate}
\end{document}
